\addtocontents{toc}{\protect\Needspace{3\baselineskip}}
\originalchapter{公开重述习题与光锥补充题}
这些题面的条件实际印在 MIT 的公开作业 PDF 中, 虽有教材编号引用, 并非只给编号的缺失题. 原题教师为 B. Zwiebach; 中文题面为重编, 下列解答均由本稿独立补写. HW4.5 与 HW5.1 是同题重发, 以一个题面与解答对应两个原始编号. 最后一题只覆盖 HW10.9 的公开补充条件, 其未取得的教材主体和 (b) 不计为完成.

\src{8.251 HW4.2, (a)--(b), 15 分; 公开作业说明为教材题 5.7 的扩充.}
\begin{exercise}
曲时空度规为 $g_{\mu\nu}(x)$, 类时粒子的正固有长度定义为 $\dd s^2=-g_{\mu\nu}\dd x^\mu\dd x^\nu$, 质量 $m>0$, 作用量 $S=-mc\int\dd s$.
(a) 从变分证明
\[
 \frac{\dd}{\dd s}\left(g_{\rho\mu}\frac{\dd x^\mu}{\dd s}\right)
 =\frac12\partial_\rho g_{\mu\nu}\frac{\dd x^\mu}{\dd s}\frac{\dd x^\nu}{\dd s}.
\]
(b) 展开左边并证明 $\dd^2x^\lambda/\dd s^2+\Gamma^\lambda_{\mu\nu}(\dd x^\mu/\dd s)(\dd x^\nu/\dd s)=0$, 其中
$\Gamma^\lambda_{\mu\nu}=\frac12g^{\lambda\rho}(\partial_\mu g_{\rho\nu}+\partial_\nu g_{\rho\mu}-\partial_\rho g_{\mu\nu})$, 且说明下标对称性及常度规极限.
\end{exercise}
\begin{solution}
不要在变分前约定所有变分路径都以固有长度作参数. 先用任意参数 $\tau$, 令 $q=-g_{\mu\nu}\dot x^\mu\dot x^\nu>0$. 有
\[
 \frac{\partial L}{\partial\dot x^\rho}=mc\frac{g_{\rho\mu}\dot x^\mu}{\sqrt q},
 \quad\frac{\partial L}{\partial x^\rho}=
 \frac{mc}{2\sqrt q}\partial_\rho g_{\mu\nu}\dot x^\mu\dot x^\nu.
\]
欧拉--拉格朗日方程两边除 $mc$, 再用 $\dd s=\sqrt q\dd\tau$ 转换, 即得 (a). 展开后
\[
 g_{\rho\mu}\frac{\dd^2x^\mu}{\dd s^2}
 +\partial_\nu g_{\rho\mu}\frac{\dd x^\nu}{\dd s}\frac{\dd x^\mu}{\dd s}
 -\frac12\partial_\rho g_{\mu\nu}\frac{\dd x^\mu}{\dd s}\frac{\dd x^\nu}{\dd s}=0.
\]
速度乘积对 $\mu,\nu$ 对称, 可将第二项替为
$\frac12(\partial_\nu g_{\rho\mu}+\partial_\mu g_{\rho\nu})u^\mu u^\nu$. 乘逆度规 $g^{\lambda\rho}$ 即得 (b). 同时交换 $\mu,\nu$ 不改变三项组合, 所以 $\Gamma^\lambda_{\mu\nu}=\Gamma^\lambda_{\nu\mu}$. 常度规在该坐标系中所有偏导为零, 恢复以固有长度仿射参数化的直线. 此练习使用 GR 卷已有度规概念, 不重复完整曲率理论.
\end{solution}

\src{8.251 HW4.5 与 HW5.1, 各 10 分, 同一题面; 两份作业均完整重述教材题 6.2.}
\begin{exercise}
相对论弦两端固定在 $(0,\boldsymbol0)$ 与 $(a,\boldsymbol0)$, $a>0$. 张力为 $T_0$. 在静态规范中取 $X^0=ct$, $X^1=x=a\sigma/\sigma_1$, 横向坐标集合为 $\boldsymbol y(t,x)$. 此参数化要求 $X^1$ 沿弦单调增加, 小振动时可满足. 在 $|\boldsymbol v_\perp|\ll c$ 且横向斜率小的条件下, 证明 Nambu--Goto 作用量除一个加性常数外变成普通非相对论横波作用量, 求张力、线密度和该常数.
\end{exercise}
\begin{solution}
直接在 $(t,x)$ 参数中计算,
$h_{tt}=-c^2+\dot{\boldsymbol y}^2$,
$h_{xx}=1+\boldsymbol y_x^2$,
$h_{tx}=\dot{\boldsymbol y}\cdot\boldsymbol y_x$. 因而
\[
 -\det h=c^2(1+\boldsymbol y_x^2)-\dot{\boldsymbol y}^2
 -\dot{\boldsymbol y}^2\boldsymbol y_x^2
 +(\dot{\boldsymbol y}\cdot\boldsymbol y_x)^2.
\]
斜率和 $|\dot{\boldsymbol y}|/c$ 为同阶小量时, 后两项为四阶. 展开平方根并乘 $-T_0/c$, 得
\[
 S=-T_0a(t_f-t_i)+\frac12\int_{t_i}^{t_f}\dd t\int_0^a\dd x
 \left[\frac{T_0}{c^2}\dot{\boldsymbol y}^2-T_0\boldsymbol y_x^2\right]+O(y^4).
\]
所以普通弦张力为 $T_0$, 线密度为 $\mu_0=T_0/c^2$, 加性常数为 $-T_0a(t_f-t_i)$, 对应静止整弦能量 $T_0a$. 若在拉格朗日量而非作用量层面说“常数”, 它是 $-T_0a$. 端点固定给 $\boldsymbol y(t,0)=\boldsymbol y(t,a)=0$, 不应从展开中误读为自由端点. 从原 $\sigma$ 参数转到 $x$ 时测度与导数的比例正好抵消, 物理答案不依任意 $\sigma_1$.
\end{solution}

\src{8.251 HW8.5, (a)--(d), 20 分; 公开 PDF 为教材题 10.7 的改进重述.}
\begin{exercise}
在 $D$ 维负时间号平直时空取 $m>0$, $F_{\mu\nu}=\partial_\mu A_\nu-\partial_\nu A_\mu$, 实标量 $\phi$, 并令
\[
 \mathcal L=-\frac14F_{\mu\nu}F^{\mu\nu}-\frac12m^2A_\mu A^\mu
 -\frac12\partial_\mu\phi\partial^\mu\phi+mA_\mu\partial^\mu\phi.
\]
(a) 给定 $\delta A_\mu=\partial_\mu\epsilon$, 求 $\delta\phi$ 使拉格朗日量不变.
(b) 求 $A_\mu$ 与 $\phi$ 的场方程.
(c) 说明可取 $\phi=0$ 的规范并简化方程.
(d) 在动量空间证明 $p^2\ne-m^2$ 时无非零解, $p^2=-m^2$ 时有 $D-1$ 个独立极化. 原 PDF 的不等号已按页面视觉核对, 避免纯文字提取漏掉斜线.
\end{exercise}
\begin{solution}
后三项为 $-\frac12(mA_\mu-\partial_\mu\phi)^2$, 因此 $\delta\phi=m\epsilon$ 使此组合严格不变, $F$ 也不变. 积分分部得
\[
 \partial_\nu F^{\nu\mu}-m^2A^\mu+m\partial^\mu\phi=0,
 \qquad\partial^2\phi-m\partial_\mu A^\mu=0.
\]
取 $\epsilon=-\phi/m$, 可使新 $\phi$ 为零, 同时 $A$ 移为 $A-\partial\phi/m$. 这是变量重排, 不丢掉原来的纵向物理自由度. 方程变为
\[
 \partial_\nu F^{\nu\mu}-m^2A^\mu=0,
 \qquad\partial_\mu A^\mu=0.
\]
对第一式求散度也给第二式, 因 $m\ne0$. 于是 $(\partial^2-m^2)A^\mu=0$. 令 $A^\mu=\varepsilon^\mu\ee^{\ii p\cdot x}$, 得 $(p^2+m^2)\varepsilon^\mu=0$ 和 $p_\mu\varepsilon^\mu=0$. 非质量壳时 $\varepsilon=0$. 在质量壳且选静止系 $p^\mu=(m,0,\ldots,0)$, 横向约束给 $\varepsilon^0=0$, 其余 $D-1$ 个独立分量可自由指定. 已取 $\phi=0$ 后没有仍可改变这些极化的连续规范参数. 若 $m=0$, 不能执行除 $m$ 的规范操作, 而是 $D-2$ 个无质量向量极化加一个独立标量, 需另行描述.
\end{solution}

\src{8.251 HW9.5, (a)--(b), 15 分, 2007 年 4 月 16 日更正版; 公开题说明改编自教材 11.7.}
\begin{exercise}
(a) 计算光锥 Lorentz 生成元 $M^{+-}$ 与粒子坐标 $x^+(\tau),x^-(\tau),x^I(\tau)$ 的对易式, 检查其预期 Lorentz 变换.
(b) 对 $M^{-I}$ 与 $x^+(\tau),x^-(\tau),x^J(\tau)$ 做同样计算, 说明结果为通常 Lorentz 变换加世界线重参数化, 并求该参数 $\lambda$. 原题提示采用厄米形式
$\delta x^\mu=\frac12\{\lambda,\partial_\tau x^\mu\}$.
\end{exercise}
\begin{solution}
为将题中引用的生成元变成可自学的明确条件, 在这里从约化自由粒子相空间建立它们. 取仿射规范 $x^+=\kappa p^+\tau$, $x^I=x_0^I+\kappa p^I\tau$, $x^-=x_0^-+\kappa p^-\tau$, 其中 $p^-=(p^Ip^I+m^2)/(2p^+)$, $p^+>0$, $\kappa>0$ 为标记尺度（原卷 Test 2 与期末公式表取 $\kappa=1/m^2$, 此时 $m>0$）, $[x_0^-,p^+]=-\ii$, $[x_0^I,p^J]=\ii\delta^{IJ}$. 这与目标时间规范 $x^+=\tau$ 只差参数重标记, 但变换补偿式会不同.

由 $M^{\mu\nu}=x^\mu p^\nu-x^\nu p^\mu$ 消去随 $\tau$ 的项并取厄米序,
\[
 M^{+-}=-\tfrac12\{x_0^-,p^+\},\qquad
 M^{-I}=x_0^-p^I-\tfrac12\{x_0^I,p^-\}.
\]
先用 $[x_0^I,f(p)]=\ii\partial_{p^I}f$ 与
$[x_0^-,f(p^+)]=-\ii\partial_{p^+}f$. 对 (a) 得
\[
 [M^{+-},x^+]=\ii x^+,\qquad
 [M^{+-},x^-]=-\ii x^-,\qquad[M^{+-},x^I]=0.
\]
这正是 $[M^{\mu\nu},x^\rho]=\ii(\eta^{\mu\rho}x^\nu-\eta^{\nu\rho}x^\mu)$ 对 $\mu=+,\nu=-$ 的结果. $p^+$ 与 $x^+$ 一起变换, 所以此仿射规范无需额外补偿.

对 (b) 可先求 $[M^{-I},p^+]=-\ii p^I$,
$[M^{-I},p^J]=-\ii\delta^{IJ}p^-$,
$[M^{-I},p^-]=0$. 相应坐标对易式为
\begin{align*}
 [M^{-I},x^+]&=-\ii\kappa\tau p^I,\\
 [M^{-I},x^-]&=\frac\ii2\left\{x_0^I,\frac{p^-}{p^+}\right\},\\
 [M^{-I},x^J]&=-\ii\delta^{IJ}x^-
 +\frac\ii2\left\{x_0^I,\frac{p^J}{p^+}\right\}.
\end{align*}
例如第二式的零模项来自
$[p^-,x_0^-]=-\ii p^-/p^+$, 随 $\tau$ 的项因 $[M^{-I},p^-]=0$ 消失. 定义实无穷小变换 $\delta x=\ii\epsilon[M^{-I},x]$. 通常矢量 Lorentz 部分为 $\delta_Lx^+=\epsilon x^I$, $\delta_Lx^-=0$, $\delta_Lx^J=\epsilon\delta^{IJ}x^-$. 上述差别统一为
\[
 \delta x^\mu-\delta_Lx^\mu
 =\frac12\{\lambda,\partial_\tau x^\mu\},
 \qquad\lambda=-\epsilon\frac{x_0^I}{\kappa p^+}.
\]
$x_0^I$ 与 $p^+$ 可交换, 故 $\lambda$ 厄米. 对 $x^+$, 补偿为 $-\epsilon x_0^I$, 正好将 $\epsilon(x_0^I+\kappa p^I\tau)$ 改为 $\epsilon\kappa p^I\tau$. 对另两坐标则用 $\partial_\tau x^-=\kappa p^-$、$\partial_\tau x^J=\kappa p^J$ 检查. 若改取 $x^+=\tau$ 固定目标时间, 补偿应重算为 $-\epsilon x^I(\tau)$; 此变化来自参数约定, 不构成 Lorentz 对称失效.
\end{solution}

\src{8.251 HW10.9, 公开补充式; 教材题 13.4 主体与引用题 8.7 未取得, 仅完成此补充.}
\begin{exercise}
对非紧致平直目标中的闭弦, $X^+=\alpha'p^+\tau$, $0\leq\sigma<2\pi$, 无横杠模依 $\tau-\sigma$, 横杠模依 $\tau+\sigma$. 证明
$L_0^\perp-\bar L_0^\perp=-\frac{p^+}{2\pi}\int_0^{2\pi}\partial_\sigma X^-\dd\sigma$, 并解释经典零模差为何消失.
\end{exercise}
\begin{solution}
约束给 $X'^-=\dot X^IX'^I/(\alpha'p^+)$. 将横向场分成 $X_R(\tau-\sigma)+X_L(\tau+\sigma)$, 则
$\dot X^I=X_R^{I\prime}+X_L^{I\prime}$,
$X'^I=-X_R^{I\prime}+X_L^{I\prime}$, 乘积为
$\dot X^IX'^I=(X_L^{I\prime})^2-(X_R^{I\prime})^2$.
两组交叉项因普通经典乘积交换而抵消. 每组导数为
$\sqrt{\alpha'/2}\sum_n\alpha_n\ee^{-\ii nu}$ 或横杠式, 利用周期 Fourier 正交性,
\[
 \frac1{2\pi\alpha'}\int_0^{2\pi}\dot X^IX'^I\dd\sigma
 =\bar L_0^\perp-L_0^\perp.
\]
代回即得所求. 非紧致 $X^-$ 严格周期使积分为 $X^-(2\pi)-X^-(0)=0$, 故经典左右层相等. 量子正规序时两边相同截距也抵消. 这没有补出未取得的原教材 (a) 其他条件和 (b), 其缺口仍列在来源表.
\end{solution}
